\section{The Cox Program, Dissected}
\label{sec:cox}

Cox's theorem~\cite{cox1946,jaynes2003} derives the probability
calculus from conditions on plausibility.  After Halpern's
counterexample~\cite{halpern1999}, its repaired
forms~\cite{vanhorn2003,arnborg-sjodin2000} fix the calculus up to
monotone regraduation.  Every known derivation is classical, and
Section~\ref{sec:hinge} showed why: the negation axiom R3 \emph{is}
excluded middle.  Dropping the R3 axiom has the following 
effects: product rule is unaffected (as it is logic-free),
the sum rule becomes an axiom rather than a result.
As a result, the naive constructive statement (strict monotonicity
on all of $\mathbb{R}$) is false.

\subsection{A Cox Structure without Negation}

\begin{lstlisting}
structure CoxModel (Ω : Type*) [Order.Frame Ω] where
  pl : Ω → Ω → ℝ            -- pl a b : plausibility of a given b
  mono_left : ∀ c, Monotone fun a => pl a c
  pl_bot : ∀ c, pl ⊥ c = 0
  pl_top : ∀ c, pl ⊤ c = 1
  F : ℝ → ℝ → ℝ             -- conjunction functional
  conj : ∀ a b c, pl (a ⊓ b) c = F (pl a (b ⊓ c)) (pl b c)
  F_assoc : ∀ x y z, F (F x y) z = F x (F y z)
  F_strictMono_left : ∀ y, 0 < y → StrictMono fun x => F x y
\end{lstlisting}

The missing R3 axiom can be observed in the omission of 
any constraint on the relation between
$\mathrm{pl}(a^{\mathsf{c}}, b)$ and $\mathrm{pl}(a, b)$.  
The $\mathtt{CoxModel}$ restricts monotonicity of $F$ to
$0 < y$ to avoid a contradiction with the fact that 
the boundary axioms force $F(x, 0) = 0$ for all
$x$.  \texttt{coxModelENNReal} (explicit model on the
non-classical chain $\ENNReal$) satisfies 
this constraint, as $a = \top$ is provably decidable.

\subsection{The Sum Rule Is Irreducible}

Classically, additivity is \emph{derived}: R3 turns disjunction into
negated conjunction via De~Morgan, and the product rule does the
rest.  Constructively that route is closed, since
$a \vee b \neq \neg(\neg a \wedge \neg b)$ in general:

\begin{theorem}[\texttt{modularity\_irreducible},
  \texttt{no\_disjunction\_\allowbreak functional}]
\label{thm:no-functional}
Let $\Omega$ be the five-element frame of lower sets of the poset
$o < a,\, o < b$ with $a, b$ incomparable.  Then there is a
monotone, normalized assignment $q$ on $\Omega$ that is additive on
disjoint joins yet not modular, and no binary operation $S$ on
values satisfies $q(x \sqcup y) = S(q(x), q(y))$ for all $x, y$.
Moreover, the
pairs $({\downarrow}a, {\downarrow}a)$ and $({\downarrow}a,
{\downarrow}b)$ have identical value-pairs $(\tfrac12, \tfrac12)$
but joins valued $\tfrac12$ and $1$.
\end{theorem}

\noindent
Unlike for a conjunction, a disjunction's plausibility cannot 
generally be expressed as its disjuncts'
plausibilities via conditioning.  So, a constructive Cox
theorem \emph{must} posit the sum-rule half.  Modularity, 
which says that the calculus is closed
under mixing certainties, is what
posits this, justified by Theorem~\ref{thm:mixture}.

\subsection{Cox Model Theorem}
 
Strict
monotonicity forces $g(-1) < g(0) = 0$, impossible in $\ENNReal$.
For this reason, it is impossible to define a strictly monotone
regraduation $g : \mathbb{R} \to \ENNReal$ with $g(0) = 0$.
The repair (strictness only on $[0,1]$, where plausibilities live)
is exactly the same domain restriction that Halpern's critique of
the classical theorem turns on.  The corrected statement is:

\begin{theorem}[\texttt{constructive\_cox}]
\label{thm:cox}
For every frame $\Omega$, every Cox model on $\Omega$ whose
unconditional plausibility is modular
regraduates to a valuation: there exist $v : \mathrm{Valuation}\;
\Omega$ and $g$ strictly monotone on $[0,1]$ with $g(0)=0$,
$g(1)=1$, and for all
$a \in \Omega$, $v(a) = g(\mathrm{pl}(a, \top))$ .
\end{theorem}

The regraduation for the sum-rule half is straightforward
($g = \texttt{ENNReal.ofReal}$).  The 
product-rule half, showing $F$ regraduates to multiplication, 
on the other hand, requires additional machinery. In particular,
we required a proof of Acz\'el's associativity
theorem~\cite{aczel1966}, a statement of
pure analysis whose type never mentions $\Omega$.
\textsf{mathlib} has no Acz\'el theorem, no H\"older-style
embedding for ordered semigroups on intervals, and no Cauchy
functional equation under monotonicity (only under continuity), so
we build the analytic core from scratch.  A
H\"older-style embedding needs two structural facts beyond
associativity and monotonicity.  The first is \emph{divisibility}:
every element must have a $2^n$-th root for each $n$, since the
generator is built as the dyadic limit $g(b) = \bigsqcup_n
(\text{count of } 2^n\text{-th roots under } b)/2^n$, which cannot
even be written down without these roots existing.  The second is
the \emph{Archimedean property}~\cite{hoelder1901}, which rules out
elements too fine for any real number to detect, and is what makes
an embedding into $\mathbb{R}$ possible at all.  
Taking these together allowed us to define a normalized, strictly monotone additive
generator and hence a strictly monotone multiplicative one on the
positive cone (\texttt{exists\_mul\_generator},
\texttt{aczelStatement\_cone}).  Our contribution shortens this
hypothesis list on both counts: the Archimedean property is
\emph{derived} from continuity and positivity, and the interval 
form of divisibility needed by the
dyadic construction is itself a consequence of the intermediate
value theorem (\texttt{exists\_diag\_eq}) applied to
\texttt{Scale}'s weaker, plain divisibility hypothesis.  The
constructive/classical divide is therefore \emph{entirely} in the
sum rule: this whole apparatus is a statement of pure analysis
whose type never mentions $\Omega$.
Appendix~\ref{app:aczel} gives the construction.  The residual
analytic gap is discussed in
Section~\ref{sec:conclusion}.
